\section{Introduction}
\label{sec:introduction}

Credit decisions determine which data become available for subsequent learning. Repayment outcomes are observed for approved applicants, while the outcomes of rejected applications remain unknown. This selective-label problem can leave evaluation and retraining dependent on the policy that collected the data \cite{lakkaraju2017selective,kleinberg2018human}. Acquiring feedback from a previously rejected region therefore requires a lending decision and acceptance of the associated credit exposure.

Existing research addresses this dependence through reject inference \cite{banasik2007reject}, decision policies for selective observations \cite{wei2021selective}, optimistic exploration \cite{pacchiano2021plot}, and adaptive bounded debiasing \cite{yang2022bounded}. Recent credit experiments also use controlled approvals to investigate feedback loops \cite{scarone2026illusion}. These approaches motivate a practical question about resource allocation: under a predicted-loss budget, does sacrificing feasible feedback quantity for a higher information score deliver sufficient downstream benefit? Non-uniform acquisition costs are already studied in active learning \cite{donmez2008proactive,tomanek2010cost}; our focus is the comparison of quantity-preserving and composition-oriented policies when the cost is an applicant-level credit-loss proxy.

We examine this question using lexicographic risk-budgeted feedback acquisition (LRBE). At each round, the first stage computes the maximum feasible count $K_t^*$ under predicted costs $\widehat c_{it}=\widehat p_{it}\mathrm{EAD}_i\mathrm{LGD}_i$. The second stage optimizes a sample-composition score subject to that count equality. A secondary score can therefore change which applicants are selected without reducing the current maximum count. Removing the equality produces a composition-first policy and exposes the number of feasible observations exchanged for a higher score. Minimum-risk, joint-objective, and randomized policies provide complementary references.

The evaluation combines three views of the trade-off. Dynamic replay on four credit datasets measures complete policy outcomes over five budget ratios. A five-policy experiment at $\rho=0.1$ compares FAVE with bootstrap-committee information gain (IG), whose score is independent of FAVE's Fisher geometry and influence approximation. Same-state audits then give every selector the same visible set, candidates, costs, and numerical budget, while a controlled one-step experiment varies exposure heterogeneity and the evaluation weight of an underrepresented region. Dynamic policies share their initial data and incoming batches but subsequently retrain on their own feedback, so later candidate sets and budgets may differ.

The experiments reveal an asymmetric trade-off. Composition-first FAVE and IG often give up many feasible observations, while their predictive estimates vary across datasets. LRBE-IG acquires 31.3, 49.5, and 40.1 more labels than QualityFirst-IG on HMEQ, LendingClub, and UCI, respectively; its Recall intervals span both directions and its UCI Brier error increases. Against MinRisk, FAVE and IG achieve similar cumulative counts while consuming more predicted risk, again with mixed predictive estimates. The controlled experiment reproduces the immediate count and risk differences, while its prespecified Brier interactions remain uncertain.

This study provides an empirical framework for evaluating risk-budgeted acquisition. The LRBE comparison separates the maximum feasible count at the current state from the score used to choose applicants, exposing the number of observations exchanged for composition. Applying the same construction to FAVE and committee IG, with randomized controls, measures the associated changes in feedback quantity, predicted-risk use, Recall, and Brier score. Same-state analyses identify the immediate selection trade-off, and the controlled experiment tests exposure heterogeneity and target-region emphasis as possible moderators of predictive return. Together, these analyses turn a broad appeal to ``informative feedback'' into a measurable comparison with minimum-risk maximum-count selection.

\section{Related Work}
\label{sec:related}

\paragraph{Learning from selectively observed outcomes.}
Work on selective labels studies the evaluation and learning problems created by decision-dependent observation \cite{lakkaraju2017selective,kleinberg2018human}. Expert consistency and stochastic fair decisions provide ways to learn or collect data beyond existing acceptance patterns \cite{dearteaga2018expert,kilbertus2020fair}. Credit reject inference examines augmentation and sample-selection correction \cite{banasik2007reject,marshall2010selection}, while finite-domain policy optimization and one-sided-feedback learning formalize sequential decisions with missing outcomes \cite{wei2021selective,jiang2021onesided}. Our replay retains observed historical outcomes and reveals them according to simulated approvals. Its estimand is acquisition-policy behavior under replay, distinct from outcome imputation and field experimentation on actual rejected borrowers.

\paragraph{Exploration objectives and costs.}
PLOT uses optimistic pseudo-labels to counter self-confirming rejection and gives regret guarantees under stated modeling assumptions \cite{pacchiano2021plot}. Adaptive bounded exploration trades debiasing against the risk of accepting otherwise rejected individuals \cite{yang2022bounded}. Partial-feedback exploration also supports subgroup coverage and fairness objectives \cite{keswani2024partial}, while cost-efficient fairness auditing accounts for assessment, processing, and potential-default costs \cite{das2025cost}. Controlled approvals in credit scoring examine the consequences of retraining under selection bias \cite{scarone2026illusion}. These works establish the surrounding problem of costly exploration. We compare how a current predicted-loss allowance is allocated between feedback quantity and composition, using randomized policies as strategy controls for the quantity constraint and selection score.
